\documentclass[11pt]{article}
\usepackage[margin=1in]{geometry}
\usepackage{amsmath}
\usepackage{booktabs}
\usepackage{hyperref}

\title{Response Process Diagnostics for fast-mlsirm}
\author{ContextualWisdomLab}
\date{2026-07-01}

\begin{document}
\maketitle

\section{Purpose}

\texttt{fast-mlsirm} needs diagnostics that distinguish item type from
response process. Dichotomous and polytomous models may use cumulative
dominance processes or ideal-point unfolding processes. A fit statistic should
not hide that distinction.

\section{Paper Basis}

The MLSIRM and MLS2PLM papers motivate posterior predictive checks over
response summaries. Those checks require posterior draws or replicated response
matrices. The current package uses a regularized JML/MAP optimizer, so runtime
diagnostics should be held-out likelihood and point-estimate residual
diagnostics until posterior samples exist.

Tay et al. (2011) compare ideal-point and dominance IRT models for both
dichotomous and polytomous data. Roberts, Donoghue, and Laughlin's generalized
graded unfolding model gives the polytomous ideal-point direction; graded and
partial-credit models give the cumulative direction.

Fox and Glas (2001) motivate multilevel IRT models for responses nested in
schools, classrooms, or other higher-level units. Bock and Zimowski's multiple
group IRT work motivates population-level comparisons across groups. The
current implementation uses those papers to define diagnostic strata, not to
claim that \texttt{fast-mlsirm} estimates hierarchical or multiple-group
parameters.

\section{Probability Contract}

Let $Y_{pi}$ be the response for person $p$ and item $i$. A diagnostic
aggregator receives category probabilities
\[
  \Pr(Y_{pi}=k) = \pi_{pik}, \qquad k=0,\dots,K_i-1.
\]
For dichotomous models, a matrix of success probabilities can be converted to
\[
  (\Pr(Y_{pi}=0), \Pr(Y_{pi}=1)) = (1-\pi_{pi}, \pi_{pi}).
\]
For polytomous cumulative or ideal-point models, the model-specific probability
generator supplies the full tensor $\pi_{pik}$.

\section{Diagnostics}

The implemented diagnostics compute observed count, log-likelihood, deviance,
Pearson chi-square, and outfit mean-square by item and by person. Category fit
summaries compute observed score, expected score, raw residual, and standardized
residual by item-category pair. Model fit reports log-likelihood, deviance,
observed count, category count, Pearson chi-square, and mean absolute category
residual.

When person-level group or cluster identifiers are supplied, the same fitted
probabilities are aggregated into multigroup and multilevel-context summaries:
\[
  G_g = \{p: \mathrm{group}_p=g\}, \qquad
  C_c = \{p: \mathrm{cluster}_p=c\}.
\]
For each stratum, \texttt{groupfit} and \texttt{clusterfit} report likelihood,
deviance, Pearson chi-square, and outfit mean-square. For item diagnostics
inside a stratum, \texttt{group\_itemfit} and \texttt{cluster\_itemfit} report
the same item-fit summaries conditional on the supplied fitted probabilities.

\section{Candidate Comparison}

External cumulative, ideal-point, dichotomous, or polytomous models can provide
candidate probability tensors. \texttt{response\_process\_dimensionality\_diagnostics}
compares those candidates with log-likelihood and deviance on the supplied
response mask. Candidate labels may represent latent dimensionality, response
process, or item type; if the probabilities are held-out predictions, the
comparison is a held-out diagnostic. The function compares the probabilities
and does not estimate the candidate models.

\begin{table}[h]
\centering
\begin{tabular}{lll}
\toprule
Scope & Current support & Deferred \\
\midrule
Dichotomous MLS2PLM & point-estimate diagnostics & posterior PPC \\
Polytomous cumulative & probability-tensor diagnostics & GRM/GPCM fitting \\
Polytomous ideal point & probability-tensor diagnostics & GGUM fitting \\
Dimensionality & K-fold held-out likelihood & Bayesian model comparison \\
Multigroup & stratum fit summaries & multiple-group estimation \\
Multilevel & cluster fit summaries & hierarchical estimation \\
\bottomrule
\end{tabular}
\end{table}

\section{Product Contract}

The API and CLI should ask users for item type, response process, observed
responses, and model probabilities before reporting fit. This keeps the
diagnostic output honest: the package can aggregate fit for supplied
probabilities now, while estimation for additional model families remains a
separate model-design task.

\end{document}
